\section{Algorithms}
\label{sec:algorithms}

We evaluate only full-coverage algorithms that produce shortest-representation
output for all IEEE~754 binary64 inputs without fallback to arbitrary
precision.  This excludes Grisu3~\cite{loitsch2010grisu}, which falls back
to multiprecision arithmetic for approximately 0.5\% of inputs,
confounding the algorithmic comparison in a specification-constrained
context where every input must be handled by the same code path.

\subsection{Definitions}

\begin{definition}[Minimal Significand]
\label{def:minimal-significand}
Given a finite IEEE~754 binary64 value $x$, the \emph{minimal
significand} is the shortest decimal significand $d$ (with digit count
$k$) such that parsing $d \times 10^{n-k}$ under roundTiesToEven
recovers $x$ exactly.
\end{definition}

\begin{definition}[Minimal Printed String]
\label{def:minimal-printed-string}
The \emph{minimal printed string} of $x$ is the output of \ecma
\texttt{Number::toString} applied to $x$: the minimal significand
formatted according to the \ecma branch table.
\end{definition}

\begin{remark}
Definitions~\ref{def:minimal-significand} and
\ref{def:minimal-printed-string} are not equivalent in string length.
Consider $x = 10^{20}$: the minimal significand is
$d = 1$ ($k = 1$, $n = 21$).  Since $n = 21 \leq 21$, \ecma uses fixed
notation, yielding \texttt{100000000000000000000} (21~characters).  Now
consider $x = 10^{21}$: again $d = 1$ ($k = 1$, $n = 22$), but $n > 21$
triggers exponential notation, yielding \texttt{1e+21} (5~characters).
The $n = 21$ boundary creates divergence between significand minimality
and printed-string length.  This means implementations cannot optimize
for shorter output strings because the \ecma formatting is spec-mandated.
\end{remark}

\subsection{Burger--Dybvig (Multiprecision)}

The \burgerdybvig algorithm~\cite{burger1996printing} uses arbitrary-precision
rational arithmetic to compute the shortest decimal representation.  It
maintains big-integer numerator/denominator pairs and iteratively
extracts decimal digits.  The algorithm is general (handles any radix)
but requires heap-allocated multiprecision integers, with cost
proportional to the magnitude of the exponent.

In our Go implementation, this uses \texttt{math/big.Int} operations.
The \texttt{sync.Pool} partially mitigates allocation cost by reusing
big-integer buffers across calls.

\subsection{Schubfach (Fixed-Width)}

The \schubfach algorithm~\cite{giulietti2022schubfach} uses fixed-width
64/128-bit arithmetic, relying on \texttt{math/bits} for 128-bit
multiplication (\texttt{Mul64}, \texttt{Add64}) and bit inspection
(\texttt{TrailingZeros64}).  It uses precomputed tables of powers of ten
and 128-bit multiplications to determine the shortest decimal
significand.  No heap allocation is required in the digit-generation
path.

The key advantage is constant-time (with respect to exponent magnitude)
operation and zero GC pressure.  The precomputed table is approximately
5~KiB.

\subsection{Ry\=u (Fixed-Width)}

The \ryualg algorithm~\cite{adams2018ryu} was the first to achieve full
coverage shortest-output generation using only fixed-width 128-bit
arithmetic and precomputed lookup tables.  Like \schubfach, it computes
a rounding interval around the input value and converts it to decimal
via table-based 128-bit multiplication.  \ryualg uses a
\texttt{lastRemovedDigit} heuristic for digit trimming: digits are
removed one at a time from both ends of the interval, tracking the last
removed digit for round-to-even tie-breaking.  Our implementation shares
the same interval computation as \schubfach (symmetric for regular
values, asymmetric at power-of-two boundaries).

\subsection{Dragonbox (Fixed-Width)}

The \dragonbox algorithm~\cite{jeon2021dragonbox} extends the
fixed-width approach with two key optimizations: a \emph{shorter
interval} case for power-of-two boundaries (where the lower half of the
rounding interval is narrower), and \emph{two-digit-at-a-time} trimming
(dividing by 100 instead of 10 in the digit extraction loop).  Our
implementation separates the shorter-interval and general cases into
distinct functions, preserving the algorithm's structural innovation.

\subsection{ECMA-262 Formatting Stage}

All four algorithms share an identical formatting stage that converts the
minimal significand $(d, k, n)$ triple into the final \ecma-compliant
string.  This stage implements the branch table from \ecma \S6.1.6.1.20
and is pure string manipulation with no arithmetic.
